\documentclass[11pt]{article}

\usepackage{amsmath, amssymb, amsfonts}
\usepackage{geometry}
\geometry{margin=1in}

\title{Supplementary Information: Information Distortion, Effective Sample Size, Corrected Covariance, and Bias}
\author{}
\date{}

\begin{document}

\maketitle

\section*{1. Log-Likelihood and Score Definitions}

Consider a trajectory of $T$ sequential measurement intervals.
The total log-likelihood is defined as
\begin{equation}
\ell(\theta) = \log L(\theta) = \sum_{t=1}^T \ell_t(\theta),
\end{equation}
with total score
\begin{equation}
S(\theta) = \nabla_\theta \ell(\theta)
= \sum_{t=1}^T s_t(\theta),
\qquad
s_t(\theta) = \nabla_\theta \ell_t(\theta).
\end{equation}

Expectations and covariances are taken over independent realizations of full trajectories generated under parameter $\theta$.

\section*{2. Gaussian Fisher Information (GFI)}

At each interval, the predictive likelihood is approximated as Gaussian:
\begin{equation}
y_t \sim \mathcal{N}\!\left(\mu_t(\theta), \sigma_t^2(\theta)\right),
\end{equation}
where $\mu_t(\theta)$ and $\sigma_t^2(\theta)$ are obtained recursively.

For scalar variance $\sigma_t^2$, the per-interval Gaussian Fisher information is
\begin{equation}
\mathbf I_t^{\mathrm{G}}(\theta)
=
\frac{1}{\sigma_t^2(\theta)}
\nabla_\theta \mu_t(\theta)
\nabla_\theta \mu_t(\theta)^\top
+
\frac{1}{2\sigma_t^4(\theta)}
\nabla_\theta \sigma_t^2(\theta)
\nabla_\theta \sigma_t^2(\theta)^\top.
\end{equation}

The total Gaussian Fisher Information is
\begin{equation}
\mathbf H(\theta)
=
\sum_{t=1}^T \mathbb{E}\!\left[\mathbf I_t^{\mathrm{G}}(\theta)\right]
=
- \mathbb{E}\!\left[\nabla_\theta^2 \ell(\theta)\right].
\end{equation}

The matrix $\mathbf H$ defines the local quadratic geometry implied by the Gaussian predictive model.

\section*{3. Score Fluctuation Matrices}

To diagnose model failure, we compare predicted curvature $\mathbf H$ against observed score fluctuations. We distinguish between local noise and trajectory-wide dependencies.

The total score covariance is
\begin{equation}
\mathbf J_{\mathrm{total}}(\theta)
=
\mathrm{Cov}\!\left(S(\theta)\right).
\end{equation}

This matrix includes both per-interval score variance and cross-interval covariances.

We also define the sample score covariance
\begin{equation}
\mathbf J_{\mathrm{sample}}
=
\sum_{t=1}^T \mathrm{Cov}(s_t).
\end{equation}

Here $\mathbf J_{\mathrm{sample}}$ is the sum of within-interval score covariances and excludes cross-interval covariances. Equivalently,
\begin{equation}
\mathbf J_{\mathrm{total}}
=
\mathbf J_{\mathrm{sample}} + 2\sum_{i<j}\mathrm{Cov}(s_i,s_j).
\end{equation}

\section*{4. Information Distortion Matrix (IDM)}

The total Information Distortion Matrix is
\begin{equation}
\mathbf C(\theta)
=
\mathbf H^{-1/2}
\mathbf J_{\mathrm{total}}
\mathbf H^{-1/2}.
\end{equation}

For full-rank $\mathbf H$, $\mathbf H^{-1/2}$ denotes the inverse square root of $\mathbf H$.
When $\mathbf H$ is singular or nearly singular, the same expression is evaluated on the retained eigenspace of $\mathbf H$ described in Section~8.
In computations, we obtain $\mathbf H^{-1/2}$ via a Cholesky factorization
$\mathbf H = \mathbf L \mathbf L^\top$ and triangular solves, without forming explicit matrix inverses.

If the predictive model is fully consistent,
\begin{equation}
\mathbf J_{\mathrm{total}} = \mathbf H,
\qquad
\mathbf C = \mathbf I.
\end{equation}

In this regime, the predictive model acts as an effective whitening transformation for the information channel: total score fluctuations match the curvature predicted by the Gaussian Fisher geometry.

\section*{5. Per-Parameter and Spectral Distortion}

In the original parameter basis, the diagonal elements of $\mathbf C$ quantify the variance inflation (or deflation) of each parameter relative to the Gaussian Fisher prediction:
\begin{equation}
C_{ii}
=
\frac{\text{real fluctuation variance in direction } i}
     {\text{Gaussian Fisher predicted variance in direction } i}.
\end{equation}

Thus:
\begin{itemize}
\item $C_{ii} = 1$ indicates consistency for parameter $i$.
\item $C_{ii} > 1$ indicates underestimation of uncertainty.
\item $C_{ii} < 1$ indicates overestimation of uncertainty.
\end{itemize}

While $\mathbf C$ contains full directional information, its diagonal provides a physically interpretable per-parameter distortion measure.

Let $\lambda_i$ denote the eigenvalues of $\mathbf C$. Then:
\begin{itemize}
\item $\lambda_i = 1$ indicates perfect agreement between curvature and fluctuation geometry along an eigen-direction.
\item $\lambda_i > 1$ indicates variance inflation relative to the Gaussian Fisher prediction.
\item $\lambda_i < 1$ indicates deflation.
\end{itemize}

The determinant $\det(\mathbf C)$ quantifies global information-volume distortion, while $\mathrm{tr}(\mathbf C)/d$ measures average directional distortion.

\section*{6. Distortion-Corrected Covariance (DCC)}

To derive the corrected covariance for parameter estimates $\hat{\theta}$, perform a first-order Taylor expansion of the score around the true parameters $\theta_0$:
\begin{equation}
S(\hat{\theta}) \approx S(\theta_0) + \nabla_\theta S(\theta_0)(\hat{\theta} - \theta_0).
\end{equation}

At the maximum likelihood estimate, $S(\hat{\theta}) = 0$. Replacing the empirical Hessian with its expectation $-\mathbf H$, we obtain
\begin{equation}
(\hat{\theta} - \theta_0) \approx \mathbf H^{-1} S(\theta_0).
\end{equation}

The covariance of the estimator is therefore approximated by the sandwich form
\begin{equation}
\Sigma_{\mathrm{DCC}}
=
\mathbf H^{-1}
\mathbf J_{\mathrm{total}}
\mathbf H^{-1}
=
\mathbf H^{-1/2}
\mathbf C
\mathbf H^{-1/2}.
\end{equation}

When $\mathbf H$ is singular, $\mathbf H^{-1}$ is understood as the inverse on the retained eigenspace of $\mathbf H$.

When $\mathbf C = \mathbf I$, $\Sigma_{\mathrm{DCC}} = \mathbf H^{-1}$.

\section*{7. Distortion-Induced Bias (DIB)}

If the expected total score is nonzero,
\begin{equation}
\mathbf g = \mathbb{E}[S(\theta)] \neq 0,
\end{equation}
then the estimator exhibits a first-order bias
\begin{equation}
\mathbf b_{\mathrm{DIB}}
=
\mathbf H^{-1} \mathbf g.
\end{equation}

In singular cases, this bias is defined only on the retained eigenspace of $\mathbf H$.

This represents the systematic parameter displacement induced by persistent score drift.

For $N$ independent trajectories,
\begin{equation}
\Sigma_{\mathrm{DCC}} \propto \frac{1}{N},
\end{equation}
while structural bias remains approximately constant. Thus the relative importance of bias increases with sample size.

\section*{8. Singular Fisher Geometry and Retained Subspaces}

The formulas above are written using inverse powers of $\mathbf H$ and $\mathbf J_{\mathrm{sample}}$.
These inverses are only literal full-space inverses when the corresponding matrices are full rank.
If $\mathbf H$ is singular, only a subset of parameter directions is identifiable under the Gaussian predictive geometry.

Let
\begin{equation}
\mathbf H = \mathbf U_H \mathbf \Lambda_H \mathbf U_H^\top
\end{equation}
be the spectral decomposition of $\mathbf H$, and retain only eigenvalues above a numerical cutoff.
The corresponding columns of $\mathbf U_H$ define the identifiable subspace of the parameter space.
All inverse powers of $\mathbf H$ are then taken only on this retained eigenspace.

Accordingly, the total distortion and sample distortion are evaluated on $\mathrm{range}(\mathbf H)$:
\begin{equation}
\mathbf C_H
=
\mathbf \Lambda_H^{-1/2}
\mathbf U_H^\top \mathbf J_{\mathrm{total}} \mathbf U_H
\mathbf \Lambda_H^{-1/2},
\end{equation}
\begin{equation}
\mathbf C_{\mathrm{sample},H}
=
\mathbf \Lambda_H^{-1/2}
\mathbf U_H^\top \mathbf J_{\mathrm{sample}} \mathbf U_H
\mathbf \Lambda_H^{-1/2}.
\end{equation}

If $\mathbf J_{\mathrm{sample}}$ is singular, the correlation distortion is defined on the retained eigenspace of $\mathbf J_{\mathrm{sample}}$ instead:
\begin{equation}
\mathbf R_{J_{\mathrm{sample}}}
=
\mathbf \Lambda_{J_{\mathrm{sample}}}^{-1/2}
\mathbf U_{J_{\mathrm{sample}}}^\top \mathbf J_{\mathrm{total}} \mathbf U_{J_{\mathrm{sample}}}
\mathbf \Lambda_{J_{\mathrm{sample}}}^{-1/2}.
\end{equation}

The corrected covariance is likewise evaluated on the retained eigenspace of $\mathbf H$:
\begin{equation}
\Sigma_{\mathrm{DCC},H}
=
\mathbf U_H
\mathbf \Lambda_H^{-1}
\mathbf U_H^\top
\mathbf J_{\mathrm{total}}
\mathbf U_H
\mathbf \Lambda_H^{-1}
\mathbf U_H^\top.
\end{equation}

If the retained eigenspace is empty, the corresponding distortion quantity is non-informative and should be treated as undefined rather than as a finite full-space matrix.

\section*{9. Effective Sample Size and Correlation Distortion}

\subsection*{8.1 Variance of Sum vs Sum of Variances}

For a sequence $X_t$, define
\begin{equation}
L_X = \sum_{t=1}^T X_t.
\end{equation}

Then
\begin{equation}
\mathrm{Var}(L_X)
=
\sum_{t=1}^T \mathrm{Var}(X_t)
+
2\sum_{i<j}\mathrm{Cov}(X_i, X_j).
\end{equation}

Define the inflation factor
\begin{equation}
\kappa_X
=
\frac{\mathrm{Var}(L_X)}
{\sum_{t=1}^T \mathrm{Var}(X_t)}.
\end{equation}

If $X_t$ are independent, $\kappa_X = 1$. The effective sample size is
\begin{equation}
T_{\mathrm{eff}}^{(X)} = \frac{T}{\kappa_X}.
\end{equation}

\subsection*{8.2 Residual-Based Effective Sample Size}

For standardized residuals
\begin{equation}
r_t = \frac{y_t - \mu_t}{\sigma_t},
\end{equation}
define
\begin{equation}
\kappa_r
=
\frac{\mathrm{Var}\!\left(\sum r_t\right)}
{\sum \mathrm{Var}(r_t)}.
\end{equation}

\subsection*{8.3 Score-Based Effective Sample Size}

Define the correlation distortion matrix
\begin{equation}
\mathbf R
=
\mathbf J_{\mathrm{sample}}^{-1/2}
\mathbf J_{\mathrm{total}}
\mathbf J_{\mathrm{sample}}^{-1/2}.
\end{equation}

If there is no temporal dependence in the score process, then $\mathbf R = \mathbf I$.

Scalar effective sample size measures may be derived from invariants of $\mathbf R$:
\begin{align}
\kappa_{\mathrm{tr}} &= \frac{1}{d}\mathrm{tr}(\mathbf R), \\
\kappa_{\mathrm{det}} &= \exp\!\left(\frac{1}{d}\log\det(\mathbf R)\right).
\end{align}

If $\mathbf R \approx \kappa \mathbf I$, temporal dependence acts isotropically.

\section*{10. Sample Distortion and Decomposition}

Define the sample distortion matrix
\begin{equation}
\mathbf C_{\mathrm{sample}}
=
\mathbf H^{-1/2}
\mathbf J_{\mathrm{sample}}
\mathbf H^{-1/2}.
\end{equation}

If $\mathbf J_{\mathrm{sample}} = \mathbf H$, then $\mathbf C_{\mathrm{sample}} = \mathbf I$.
Thus $\mathbf C_{\mathrm{sample}}$ measures mismatch between within-interval score variance and Gaussian Fisher curvature.

The total Information Distortion Matrix decomposes as
\begin{equation}
\mathbf C
=
\mathbf C_{\mathrm{sample}}^{1/2}
\mathbf R
\mathbf C_{\mathrm{sample}}^{1/2},
\end{equation}
with
\begin{equation}
\mathbf R
=
\mathbf C_{\mathrm{sample}}^{-1/2}
\mathbf C
\mathbf C_{\mathrm{sample}}^{-1/2}
\end{equation}
in Fisher coordinates.

If distortion is purely due to temporal dependence and acts isotropically, then
\begin{equation}
\mathbf C \approx \kappa \mathbf I
\quad \Rightarrow \quad
\Sigma_{\mathrm{DCC}} \approx \kappa \mathbf H^{-1}.
\end{equation}

In this regime, a simple effective sample size correction captures most of the distortion.
However, if $\mathbf R$ is anisotropic or $\mathbf C_{\mathrm{sample}} \neq \mathbf I$, then distortion cannot be reduced to scalar sample-size inflation.

\section*{11. Summary}

\begin{align}
\mathbf H &\rightarrow \text{Gaussian Fisher curvature}, \\
\mathbf J_{\mathrm{total}} &\rightarrow \text{total score fluctuation geometry}, \\
\mathbf J_{\mathrm{sample}} &\rightarrow \text{within-interval score fluctuation geometry}, \\
\mathbf C &\rightarrow \text{total information distortion}, \\
\mathbf C_{\mathrm{sample}} &\rightarrow \text{sample distortion}, \\
\mathbf R &\rightarrow \text{correlation distortion}, \\
\Sigma_{\mathrm{DCC}} &\rightarrow \text{distortion-corrected covariance}, \\
\mathbf b_{\mathrm{DIB}} &\rightarrow \text{distortion-induced bias}.
\end{align}

Together these quantities distinguish temporal dependence, sample-level mismatch, variance inflation, and systematic bias.

\end{document}
